% ---------------------------------------------------------------------
%  CHƯƠNG 1 — MỞ ĐẦU
% ---------------------------------------------------------------------
\chapter{Mở đầu}
\label{ch:modau}

\begin{muctieu}
Chương mở đầu trình bày lý do chọn đề tài, mục tiêu cần đạt, phạm vi thực
hiện và bố cục báo cáo.
\end{muctieu}

\section{Lý do chọn đề tài}

Trong một giải đấu thể thao dạng vòng tròn, kết quả các trận được ghi nhận rải
rác tại nhiều trận đấu diễn ra cùng lúc. Mỗi khi có một trận kết thúc, bảng
xếp hạng phải được tính lại và phát tới mọi người đang theo dõi, thường là trên
nhiều thiết bị khác nhau. Nếu làm thủ công thì chậm và dễ nhầm; nếu để nhiều
người cùng ghi vào một bảng chung mà không đồng bộ thì kết quả có thể bị mất
hoặc trùng. Bài toán này phù hợp với mô hình máy khách--máy chủ: một máy chủ
giữ trạng thái chung, các trọng tài là máy khách gửi kết quả, các khán giả là
máy khách nhận bảng xếp hạng cập nhật. Đề tài được chọn vì nó gom được nhiều
chủ đề của học phần vào cùng một ứng dụng vừa và nhỏ: socket TCP đa luồng
\en{(chapter 4)}, cấu trúc dữ liệu và streams \en{(chapter 2, 3)} và các tính
năng hiện đại của Java \en{(chapter 6)}.

\section{Mục tiêu của đồ án}

Đồ án hướng tới các mục tiêu sau:
\begin{itemize}[leftmargin=1.6em, itemsep=2pt]
  \item Xây dựng được máy chủ TCP đa luồng tiếp nhận kết quả và phát bảng xếp
        hạng cập nhật tới mọi khách đang xem.
  \item Đảm bảo đúng khi nhiều trọng tài gửi kết quả cùng lúc: không mất trận,
        không trùng trận, bảng xếp hạng tính lại an toàn luồng.
  \item Cài đặt đủ sáu lệnh giao thức \code{WATCH}, \code{SUBMIT},
        \code{STANDINGS}, \code{HISTORY}, \code{EXPORT}, \code{QUIT}.
  \item Tính bảng xếp hạng theo điểm, hiệu số, bàn thắng rồi phá đồng hạng bằng
        kết quả đối đầu.
  \item Lưu và khôi phục dữ liệu bằng CSV, đóng gói được thành file jar.
\end{itemize}

\section{Phạm vi và đối tượng}

Đồ án tập trung vào giải đấu vòng tròn một lượt, mỗi đội gặp mỗi đội khác đúng
một lần trên bảng xếp hạng tổng hợp. Giao thức là văn bản một dòng, tên đội
không chứa khoảng trắng. Ứng dụng chỉ dùng thư viện chuẩn của Java JDK 17 trở
lên, không dùng cơ sở dữ liệu, không dùng Spring, không bắt buộc giao diện đồ
hoạ. Đồ án không xử lý lịch thi đấu tự động ở phần bắt buộc (để ở phần mở
rộng), không hỗ trợ UDP và không phân tán máy chủ.

\begin{luuy}
Phạm vi được thu hẹp đúng mức một ứng dụng vừa và nhỏ: lớn hơn một bài tập
chương nhưng nhỏ hơn một đồ án, đúng yêu cầu của chuẩn đầu ra CLO3.
\end{luuy}

\section{Bố cục báo cáo}

Phần còn lại của báo cáo được tổ chức như sau. Chương~\ref{ch:coso} trình bày
cơ sở lý thuyết và công nghệ sử dụng. Chương~\ref{ch:thietke} phân tích bài
toán và thiết kế giải pháp. Chương~\ref{ch:caidat} mô tả quá trình cài đặt và
thử nghiệm. Chương~\ref{ch:ketqua} trình bày kết quả và đánh giá.
Chương~\ref{ch:ketluan} tổng kết và nêu hướng phát triển.
